% !TEX root = main.tex

\section{Data and Methodology}
\label{sec:methodology}

\subsection{Sample and splits}
\label{sec:data}

The sample is a balanced daily panel of eight large-capitalisation United States
technology equities---AAPL, MSFT, AMZN, GOOGL, META, NVDA, TSLA and ORCL---from
2 March 2016 to 29 December 2025. Adjusted closing prices are obtained from
Yahoo Finance and macro-financial series from the Federal Reserve Economic Data
service: the effective federal funds rate (DFF), the ten-year Treasury constant
maturity yield (DGS10), the S\&P~500 index (SP500) and the CBOE volatility index
(VIX). The market proxy is the SPY exchange-traded fund.

The forecast target is the one-day-ahead log return. For ticker $i$, a panel row
is keyed by date $t$ and its target is
\begin{equation}
  y_{i,t} \;=\; \log P_{i,t+1} - \log P_{i,t},
  \label{eq:target}
\end{equation}
realised on $t+1$. Every predictor is observable at the close of $t$, so no
input uses information from the period being forecast.

The panel is divided chronologically into three blocks, with no shuffling at any
stage:

\begin{center}
\begin{tabular}{llr}
\toprule
Split & Dates & Rows \\
\midrule
Training   & 2016-03-02 to 2019-12-30 & 7\,720 \\
Validation & 2019-12-31 to 2022-12-29 & 6\,048 \\
Test       & 2022-12-30 to 2025-12-29 & 6\,008 \\
\bottomrule
\end{tabular}
\end{center}

The split is deliberately demanding. The training block spans a low-rate,
low-volatility regime; the validation block contains the COVID-19 dislocation;
and the test block spans the 2022--2023 tightening cycle and its aftermath. Any
relation that survives is required to survive a change in monetary regime.

\subsection{Feature dictionary and stationarity}
\label{sec:features}

Model inputs are declared explicitly rather than inferred from column type. This
is a substantive design choice, not bookkeeping. Selecting every numeric column
admits level series---an index level, a policy rate---whose mean moves
permanently between the training and test blocks. Standardised on training
moments and applied to the test block, such a series arrives at a magnitude the
estimator has never seen, and any model with an unbounded linear output
extrapolates along it. Section~\ref{sec:pitfall} reports the consequences
directly.

Membership is therefore fixed in advance, and each candidate is classified by an
augmented Dickey--Fuller test. Three series are non-stationary over the sample:
the S\&P~500 level ($p=0.99$), the effective federal funds rate ($p=0.83$) and
the ten-year yield ($p=0.80$). All three are excluded and replaced by their
first differences or by a stationary counterpart. The retained set comprises
fourteen predictors: six own-return lags, three trailing realised volatilities,
the market return, the log VIX and its change, and the daily changes in the
federal funds rate and the ten-year yield. Hybrid specifications add one further
input, the GARCH conditional standard deviation of
Section~\ref{sec:garch}, so base and hybrid models differ by exactly one column.
Table~\ref{tab:features} lists every candidate, retained or excluded, with the
reason for exclusion.

Two redundancies are also removed. One column duplicated another exactly, and
the S\&P~500 return correlated $0.9985$ with the market proxy return; retaining
both would inflate collinearity and render the Bayesian coefficient posteriors
uninterpretable.

\subsection{Evaluation protocols}
\label{sec:protocols}

Two protocols are used throughout, and the difference between them is itself a
result.

Under the \emph{fixed-origin} protocol a model is estimated once on the training
block, and the fitted model together with its feature scaler is frozen and
applied to the whole test block. This is the conventional design and is
appropriate when predictors are stationary.

Under the \emph{expanding-window} protocol the model walks forward month by
month. At each fold it is refitted on every observation whose target date falls
strictly before the fold start, and it forecasts that month only; the scaler is
refitted inside the fold along with the model. This is what a forecaster could
have executed in real time.

Both protocols score the identical 6\,008 test observations, so their metrics
are directly comparable, and across all 36 folds no training observation reaches
into the month being scored.

\subsection{Conditional volatility}
\label{sec:garch}

Three conditional volatility specifications are estimated per ticker on returns
in percent: GARCH(1,1), GJR-GARCH(1,1) and EGARCH(1,1), each with Student-$t$
innovations. Parameters are estimated on the training block alone. They are then
held fixed and the recursion is filtered forward through the validation and test
blocks, so the stored quantity is a genuine one-step-ahead forecast
$\sigma_{i,t+1\mid t}$ that uses only information through $t$. The value recorded
at date $t$ therefore forecasts the same period as the target
in~\eqref{eq:target}.

Every fit is checked for optimiser convergence, finite parameters and a
conditional standard deviation within a plausible range before it is written;
the pipeline aborts otherwise. All twenty-four fits (eight tickers, three
specifications) satisfy these conditions, with conditional standard deviations
between $0.6\%$ and $13.8\%$ per day against a realised return standard
deviation of $2.5\%$. Per-fit parameters, log-likelihoods and AIC values appear
in Appendix Table~\ref{tab:garchfits}.

\subsection{Hierarchical Bayesian model}
\label{sec:bayes}

The central specification pools information across tickers while treating the
GARCH forecast as a known observation-level scale. For ticker $i$ on day $t$,
with $\mathbf{x}_{i,t}$ the standardised predictor vector,
\begin{align}
  y_{i,t} \mid \alpha_i, \boldsymbol\beta, s^2
    &\sim t_\nu\!\left(\alpha_i + \mathbf{x}_{i,t}^{\top}\boldsymbol\beta,\;
                        s^2\sigma_{i,t+1\mid t}^{2}\right),
    \label{eq:likelihood}\\
  \alpha_i \mid \mu_\alpha, \tau_\alpha^2 &\sim \mathcal{N}(\mu_\alpha, \tau_\alpha^2),
    \label{eq:pooling}\\
  \beta_j \mid \lambda_j, \tau_\beta &\sim \mathcal{N}(0, \lambda_j^2\tau_\beta^2),
  \qquad \lambda_j \sim \mathcal{C}^{+}(0,1), \quad \tau_\beta \sim \mathcal{C}^{+}(0,1).
    \label{eq:horseshoe}
\end{align}

Three features of this specification carry the analysis.

The GARCH forecast enters~\eqref{eq:likelihood} as a known scale, so the model
does not relearn volatility clustering and $s^2$ is a single scalar measuring
whether that scale is correctly sized overall. A value near one means the
volatility model is already calibrated.

Equation~\eqref{eq:pooling} is the partial-pooling layer. Because $\alpha_i$ is
drawn from a population distribution, a ticker absent from the training sample
can still be forecast by integrating $\alpha$ over that distribution, which adds
$\tau_\alpha^2$ to the predictive variance. Section~\ref{sec:transfer} exploits
this. Following \citet{gelman2006prior}, $\tau_\alpha$ receives a half-Cauchy
prior rather than a conventional inverse-gamma: with only eight groups the
between-ticker standard deviation is weakly identified, and a diffuse
inverse-gamma prior largely reproduces its own mode.

Equation~\eqref{eq:horseshoe} is the horseshoe prior of
\citet{carvalho2010horseshoe}. It replaces the more common practice of
pre-screening predictors by their correlation with the target. That practice is
unsafe in this setting for a specific reason: a trending series correlates with
almost any target over a multi-year window, so a correlation screen
preferentially selects exactly the non-stationary levels that
Section~\ref{sec:features} excludes. The horseshoe retains every predictor and
shrinks weak coefficients continuously toward zero while leaving strong ones
nearly unshrunk, so selection occurs inside the posterior and carries
uncertainty. It is the continuous analogue of the spike-and-slab selection
of \citet{qiu2018multivariate}.

Student-$t$ innovations are represented as a scale mixture of normals,
$y \sim \mathcal{N}(\mu, s^2\sigma^2/\omega)$ with
$\omega \sim \mathrm{Gamma}(\nu/2,\nu/2)$, which recovers~\eqref{eq:likelihood}
marginally and keeps every conditional in closed form.

Inference is by blocked Gibbs sampling. All conditionals are conjugate,
including the horseshoe scales, via the auxiliary-variable representation of
\citet{makalic2016simple}; the degrees of freedom $\nu$ are drawn by griddy
Gibbs. Four chains are run from dispersed starting points for 40\,000 iterations,
discarding the first 20\,000 and thinning by four, giving 20\,000 retained draws.
Convergence is assessed by rank-normalised $\widehat{R}$ and effective sample
size \citep{vehtari2021rank}: across 27 monitored quantities the largest
$\widehat{R}$ is 1.0033 and the smallest bulk effective sample size is 1\,703.

The sampler is validated on synthetic data in which three of twelve coefficients
are non-zero. It recovers all three signal coefficients with credible intervals
excluding zero, produces no false positives among the nine null coefficients,
and recovers the observation scale and the between-group standard deviation.

\subsection{Predictive distributions and scale calibration}
\label{sec:calibration}

Predictive means and scales use the full posterior rather than plug-in point
estimates, so parameter uncertainty enters the interval through the law of total
variance.

One calibration step is applied. GARCH parameters are frozen at the end of the
training block, and the training block is calmer than what follows: measured as
the standard deviation of $y/\sigma$, the scale the data require is $1.002$ on
the training block but $1.118$ on validation and $1.117$ on test. A scale fitted
to training data alone therefore understates later volatility by roughly twelve
percent. A single scalar is estimated on the \emph{validation} block and applied
to the test block; no test observation informs it.

The matching statistic matters. Because the validation block contains the
COVID-19 dislocation, a variance match is dominated by a few extreme days, which
is behaviour the Student-$t$ tails are meant to absorb rather than the scale.
Matching the interquartile range targets the body of the distribution, which is
what the scale parameter governs, and yields a factor of $1.195$ against $1.220$
for the variance match.

\subsection{Scoring rules and inference}
\label{sec:scoring}

Point accuracy is reported as root mean squared error, mean absolute error,
directional accuracy and out-of-sample $R^2$ against a stated reference
forecast. The reference is the per-ticker mean computed under the same protocol
as the model being scored, so it never enjoys an information advantage.

Probabilistic accuracy is reported as the mean logarithmic score, the continuous
ranked probability score and empirical coverage at the 50\%, 80\%, 90\% and 95\%
levels \citep{gneiting2007strictly}. All three are reported because they can
disagree: a model may be better calibrated while being less sharp. Calibration
is assessed additionally by the probability integral transform, whose values are
uniform under a correctly specified predictive distribution, tested formally by
a Kolmogorov--Smirnov statistic.

Volatility forecasts are compared under QLIKE against the squared return as the
realised-variance proxy. QLIKE is chosen because it preserves the ranking a
perfect proxy would give even when the available proxy is noisy
\citep{patton2011volatility}.

All comparisons are tested rather than asserted. Equal predictive accuracy is
assessed by the \citet{diebold1995comparing} test on per-observation loss
differentials, using a Newey--West heteroskedasticity- and
autocorrelation-consistent variance because daily loss differentials are
serially correlated, together with the small-sample correction of
\citet{harvey1997testing}. Where a family of tests is reported---one per
model--ticker pair---significance is additionally assessed after
Benjamini--Hochberg and Bonferroni correction, since at conventional levels a
handful of rejections among dozens of tests is what chance alone produces.
